\subsection{Limitations and Future Work}
\label{cic:sec:limitations}
Our methodology used surveys and interviews to elicit responses from participants regarding their past experiences and current perceptions. 

Future research should look to expand on this work through studies that focus on clear categories of smart home owners, based on characteristics such as location (e.g., temperate zones vs. frigid vs. coastal), size of home (apartment vs. single-family home vs. mansion), and other factors. In addition, our survey respondents and interview participants likely skewed more towards a tech-savvy population, especially given that we recruited using Prolific~\cite{abrokwa2021comparing}. Thus, future work can also explore perceptions and experiences of users who may be overall less aware of smart home device capabilities. \revii{Relatedly, future studies should differentiate between generations of users, who may hold different expectations of control, adopt different troubleshooting strategies, and prefer different points on the control spectrum. Older adults in particular have distinct needs: prior work built and deployed a personal assistance system that delivered medication and appointment reminders alongside wearable health sensors for independent and assisted living~\cite{hou2007pas}, and the accompanying field study found that the social structures surrounding older users shape both what such a system must do and whether it is accepted at all~\cite{alshebli2007assisted}. Because the needs and wants of each population---and indeed each country---differ substantially, these are hard and open problems that require specialized study of the human \emph{and} the systems aspects together. This thesis provides groundwork for both: the control spectrum characterized here, and the abstractions that support movement along it, give such studies a common vocabulary and a set of mechanisms to build on.}
In addition, surveys and interviews rely on participants to reflect on their interactions with smart home devices, rather than having the research team observe these interactions directly. Future work can utilize more observational methodologies and ethnographic studies to expand our results by providing additional context on how participants control their devices. 

Our results are tied to the technology mix used during data collection. However, as technology evolves and matures, and new standards such as Matter~\cite{matter} emerge, future work can incorporate our broader implications and findings into its design. Interaction patterns related to user transitions across the control spectrum will remain even as such standards develop, which will require future research to use our approach to understand \textit{how} users' interaction patterns change with these standards in place.

Our work focused primarily on a single user in a smart home; while participants did discuss how other users in the home influenced their perceptions and preferences, our study elicited only one opinion from each smart home. Previous work has described smart homes as more sociotechnical systems, often containing multiple users with diverse experiences, desires, and perceptions~\cite{geeng2019s}. Future work should more explicitly consider the sociotechnical nature by conducting studies with multiple users in the same home, examining how control schemes and methods in a single household affect different users in a smart home, and building upon the perspectives elicited in this work. 

